\documentclass[14pt, a4paper]{book}
%=========================================CÁC GÓI VÀ HÀM CHUẨN BỊ==========================================
% Cho phép sử dụng màu và vẽ viền
\usepackage{xcolor}
\usepackage{tikz}
\usepackage{contour}
\contourlength{2pt}
\usepackage{gradient-text}

\definecolor{vol}{HTML}{800000}
\definecolor{vols}{HTML}{808080}
\definecolor{eptype}{HTML}{af0158}      
\definecolor{oldpaper}{HTML}{362C03}       
\definecolor{epattr}{HTML}{f5a5d5}
\definecolor{errorone}{HTML}{92d050}
\definecolor{errortwo}{HTML}{00fff2}
\definecolor{errorthree}{HTML}{f6ff00}
\definecolor{errorfour}{HTML}{ea1712}
\definecolor{errorfive}{HTML}{ff17ff}

% Công thức toán học
\usepackage{amsmath}
\usepackage{amssymb}
\usepackage{unicode-math}
\usepackage{gensymb}

% Hộp văn bản
\usepackage[most]{tcolorbox}
\usepackage{soul}

% Liên kết
\usepackage{hyperref}

% Chèn các file nhỏ
\usepackage{subfiles}

% Chèn ảnh, video và GIF
\usepackage{graphicx}   
\usepackage{adjustbox}
\usepackage{wrapfig}
%%%%%%%%%%%%%%%%%%%%%%%%%%%%%%%%%%%%%%%%%%%%%%%%%%%%%%%%%%%%%%%%%%%%%%%%%%%%%%
% \embedvideo{<poster or text>}{<video file (MP4+H264)>}
% \embedvideo*{...}{...}                     % auto-play
%%%%%%%%%%%%%%%%%%%%%%%%%%%%%%%%%%%%%%%%%%%%%%%%%%%%%%%%%%%%%%%%%%%%%%%%%%%%%%

\usepackage[bigfiles]{pdfbase}
\ExplSyntaxOn
\NewDocumentCommand\embedvideo{smm}{
  \group_begin:
  \leavevmode
  \tl_if_exist:cTF{file_\file_mdfive_hash:n{#3}}{
    \tl_set_eq:Nc\video{file_\file_mdfive_hash:n{#3}}
  }{
    \IfFileExists{#3}{}{\GenericError{}{File~`#3'~not~found}{}{}}
    \pbs_pdfobj:nnn{}{fstream}{{}{#3}}
    \pbs_pdfobj:nnn{}{dict}{
      /Type/Filespec/F~(#3)/UF~(#3)
      /EF~<</F~\pbs_pdflastobj:>>
    }
    \tl_set:Nx\video{\pbs_pdflastobj:}
    \tl_gset_eq:cN{file_\file_mdfive_hash:n{#3}}\video
  }
  %
  \pbs_pdfobj:nnn{}{dict}{
    /Type/RichMediaInstance/Subtype/Video
    /Asset~\video
    /Params~<</FlashVars (
      source=#3&
      skin=SkinOverAllNoFullNoCaption.swf&
      skinAutoHide=true&
      skinBackgroundColor=0x5F5F5F&
      skinBackgroundAlpha=0.75
    )>>
  }
  %
  \pbs_pdfobj:nnn{}{dict}{
    /Type/RichMediaConfiguration/Subtype/Video
    /Instances~[\pbs_pdflastobj:]
  }
  %
  \pbs_pdfobj:nnn{}{dict}{
    /Type/RichMediaContent
    /Assets~<<
      /Names~[(#3)~\video]
    >>
    /Configurations~[\pbs_pdflastobj:]
  }
  \tl_set:Nx\rmcontent{\pbs_pdflastobj:}
  %
  \pbs_pdfobj:nnn{}{dict}{
    /Activation~<<
      /Condition/\IfBooleanTF{#1}{PV}{XA}
      /Presentation~<</Style/Embedded>>
    >>
    /Deactivation~<</Condition/PI>>
  }
  %
  \hbox_set:Nn\l_tmpa_box{#2}
  \tl_set:Nx\l_box_wd_tl{\dim_use:N\box_wd:N\l_tmpa_box}
  \tl_set:Nx\l_box_ht_tl{\dim_use:N\box_ht:N\l_tmpa_box}
  \tl_set:Nx\l_box_dp_tl{\dim_use:N\box_dp:N\l_tmpa_box}
  \pbs_pdfxform:nnnnn{1}{1}{}{}{\l_tmpa_box}
  %
  \pbs_pdfannot:nnnn{\l_box_wd_tl}{\l_box_ht_tl}{\l_box_dp_tl}{
    /Subtype/RichMedia
    /BS~<</W~0/S/S>>
    /Contents~(embedded~video~file:#3)
    /NM~(rma:#3)
    /AP~<</N~\pbs_pdflastxform:>>
    /RichMediaSettings~\pbs_pdflastobj:
    /RichMediaContent~\rmcontent
  }
  \phantom{#2}
  \group_end:
}
\ExplSyntaxOff
%%%%%%%%%%%%%%%%%%%%%%%%%%%%%%%%%%%%%%%%%%%%%%%%%%%%%%%%%%%%%%%%%%%%%%%%%%%%%%	
\usepackage{animate}
\usepackage[permil]{overpic}

% Lề và cỡ chữ tiêu chuẩn
\usepackage{extsizes}                       % Cho cỡ chữ 14pt
\usepackage[margin=1.54cm]{geometry}        % Căn chỉnh lề giấy
\usepackage{parskip}
\usepackage{emptypage}

% Không thụt đầu dòng khi xuống dòng mới 
\parindent 0pt

% Gạch chân xuống dòng
\usepackage{ulem}

% Thay đổi thiết lập chú thích
\usepackage{caption}
\usepackage{subcaption}

\captionsetup{justification= centering, font=small, labelformat=empty}

% Sử dụng để thay đổi font
\usepackage{fontspec}
\usepackage{luatexja}
\usepackage{luatexja-fontspec}

% Viết furigana
\usepackage{ruby}
\renewcommand{\rubysize}{0.4}
\renewcommand{\rubysep}{1pt}

% Chữ cầu vồng
\usepackage{chickenize}
% play with this value, the step is the reciprocal of the number of letters
% needed to get from red to yellow.
\chickenizesetup{rainbow_step=0.15}

% Phông cho tên-số các tập
\newfontface\flaregothic[Path=../../../Fontface/]{flaregothicbold.ttf}
\newfontface\cafeteria[Path=../../../Fontface/]{NHC-Cafeteria-Bold.ttf}
\newfontface\fotskip[Path=../../../Fontface/]{FOT-Skip Std B.otf}
\newfontface\vnmsans[Path=../../../Fontface/]{VNM Sans Std.otf}
\newfontface\vnmdisp[Path=../../../Fontface/]{VNM Sans Display Bold.otf}
\newfontface\sen[Path=../../../Fontface/]{NHC-SenBold.ttf}
\newfontface\continuum[Path=../../../Fontface/]{contb.ttf}
\newfontface\japfont[Path=../../../Fontface/]{msgothic.ttc}
\newfontface\jiga[Path=../../../Fontface/]{Jigmo.ttf}
\newfontface\jigb[Path=../../../Fontface/]{Jigmo3.ttf}
\newfontface\jigc[Path=../../../Fontface/]{Jigmo3.ttf}
\newfontface\quincy[Path=../../../Fontface/]{Quincy Caps Regular.ttf}

% Phông cho tiêu đề cuốn truyện và đầu-chân trang
\newfontface\frank[Path=../../../Fontface/]{Franklin Gothic Heavy Regular.ttf}
\newfontface\caviar[Path=../../../Fontface/]{VNF-Cariar Dreams Regular.ttf}
\newfontface\hbrk[Path=../../../Fontface/]{HouseBrokenRough.ttf}
\newfontface\ab[Path=../../../Fontface/]{angrybirds-regular.ttf}
\newfontface\abv[Path=../../../Fontface/]{0259-LNTH-angrybirds-regular.ttf}
\newfontface\device[Path=../../../Fontface/]{Device Regular.ttf}
\newfontface\harosbaelz[Path=../../../Fontface/]{BAEFONT_NORMAL-REGULAR_V1.TTF}
\newfontface\brian[Path=../../../Fontface/]{Brianne_s_hand.ttf}
\newfontface\nexa[Path=../../../Fontface/]{Nexa-Heavy.ttf}
\newfontface\digi[Path=../../../Fontface/]{DS-DIGI.TTF}
\newfontface\freshbot[Path=../../../Fontface/]{FRESHBOT.TTF}
\newfontface\arial[Path=../../../Fontface/]{arial.ttf}

% Đặt font chính và font tiếng Nhật
\setmainfont{Times New Roman}
\setmainjfont{FOT-Skip Std B}
\setmathfont
[    Extension = .otf,
     BoldFont = XITSMath-Bold,
]
{XITSMath-Regular}

% Điểu chỉnh đầu trang và chân trang
\usepackage{fancyhdr}
\usepackage{spverbatim}

% Dùng cho điều chỉnh kích thước và vị trí bảng
\usepackage{array}
\usepackage{changepage}

% Nhóm cột
\usepackage{multirow}
\usepackage{float}
% \restylefloat{figure}

% Dùng cho đoạn hội thoại
\usepackage{setspace} % Chia khoảng cách
\usepackage{dialogue}
\usepackage{calc} % Thiết lập khoảng cách giữa tên nhân vật và văn bản

% Tăng chỉ số chiều sâu của sách
\setcounter{tocdepth}{4}
\setcounter{secnumdepth}{4}

%==================================================MACRO===================================================
% A. Đường thẳng, ký hiệu
\newcommand{\ct}
{\begin{center}[---]\end{center}} 		  	                          % Cắt cảnh dài
\newcommand{\chrint}        
{\begin{center}||---||\end{center}}						        	            % Cắt đến giói thiệu nhân vật, có thể bỏ qua
\newcommand{\pov}[1]
{\begin{center}$\langle${#1}$\rangle$\end{center}}			            % Đổi góc nhìn nhân vật (mặc định là ngôi thứ ba)
\newcommand{\il}{\noindent\makebox[\textwidth]
{\rule{\paperwidth}{0.4pt}}}					  			                      % Đường thẳng phân chia

% B. Chèn ảnh và gif
\newcommand{\img}[2][12.5cm]{%
    \begin{figure}[H]
        \centering
        \includegraphics[width=#1]{../../../Image/BookER/#2.png}                    % Tự động thêm "Image/" vào trước tên tệp; chỉ cần chèn tên tệp
    \end{figure}
}

\newcommand{\gif}[4][12.5cm]{%
    \begin{figure}[H]
        \centering
        \animategraphics[autoplay,loop,width=#1]
        {30}{../../../Image/BookER/#2/#2.}{#3}{#4}                                  % Chỉ cần chèn tên tệp, lấy từ khung #3 tới #4
    \end{figure}
}

\newcommand{\pictdum}{										                          % Ảnh chèn vào thay thế ảnh gif để canh dòng
        \includegraphics[width=8cm]{placeholder.png}
}
    
\newcommand{\vid}[3][12.5cm]{%
  \embedvideo{\includegraphics[width=#1]{../../../Image/BookER/#2.png}}{#3.mp4}
}

\newcommand{\emp}[1]{%
    \includegraphics[width=2cm]{../../../Character/#1.png}
}

% C. Viết cụm và định dạng
\newcommand{\omk}{\begin{center}{\abv\Large Truyện phụ}\end{center}}
\newcommand{\anni}{\color{bdaya}A\color{bdayb}N\color{bdaya}N\color{bdayb}I\color{bdaya}V\color{bdayb}E\color{bdaya}R\color{bdayb}S\color{bdaya}A\color{bdayb}I\color{bdaya}R\color{bdayb}E}
\newcommand{\xmas}{\color{green}N\color{gray}O\color{red}Ë\color{green}L}
\newcommand{\hlw}{\color{purple}H\color{brown}A\color{purple}L\color{brown}L\color{purple}O\color{brown}W\color{purple}E\color{brown}E\color{purple}N}
\newcommand{\trva}{\begin{center}{\abv\Large Bạn có biết?}\end{center}}
\newcommand{\hll}{\textit{hololive}}
\newcommand{\stpt}[1]{\vspace{-30pt}\begin{center}{\sen\Large {#1}} \end{center}}     % Subsection
\newcommand{\stcap}[1]{\begin{center}{\caviar {#1}} \end{center}}       % Phụ đề cho subsection
\newcommand{\xdo}{\\[-24pt]}                                            % Dùng cho phần dialogue

% D. Bảng xếp hàng (cho các tập Top 10)
\newcommand{\topten}{\ruby{大}{だい}\ruby{惨}{さん}\ruby{事}{じ}}
\newcommand{\bxh}[2]{
    \begin{center}
        {{\cafeteria\Large Mystery #1:} \\
        {\textbf{#2}}}
    \end{center}
}

\newcommand{\alt}[3]{
    \begin{center}
        {(Từ {\flaregothic\textbf{ÉPISODE #1}}, quyển số #2, tại vũ trụ mang số hiệu #3.)}
    \end{center}
}

% E. Ngôn ngữ
\newcommand{\vie}[1]{\color{red} #1\color{black}}        % Tiếng Việt (Etsunan ở Round 1 và 2)
\newcommand{\eng}[1]{\color{blue} #1\color{black}}       % Tiếng Anh
\newcommand{\ind}[1]{\color{green} #1\color{black}}      % Tiếng Indonesia
\newcommand{\chn}[1]{\color{orange} #1\color{black}}     % Tiếng Trung

% F. Hộp chữ
\newcommand{\txtbx}[2][0.5]{
    \begin{tcolorbox}[
        colframe=vol,  % Màu viền
        colback=white,   % Màu nền
        boxrule=#1pt,   % Độ dày viền
        arc=0mm,         % Góc vuông (không bo tròn)
        left=5pt,        % Khoảng cách giữa nội dung và viền trái
        right=5pt,       % Khoảng cách giữa nội dung và viền phải
        top=5pt,         % Khoảng cách giữa nội dung và viền trên
        bottom=5pt       % Khoảng cách giữa nội dung và viền dưới
    ]
    #2
    \end{tcolorbox}
}

% G. Chỉnh ảnh
\newcommand{\dimtorightedge}{%
  \dimexpr\paperwidth-1in-\hoffset-\oddsidemargin\relax}

\newcommand{\dimtotop}{%
  \dimexpr\height-1in-\voffset-\topmargin-\headheight-\headsep\relax}


% Thiết đặt môi trường cho đoạn hội thoại
% A. Dùng để viết hoa chữ đầu tiên tên nhân vật
\renewcommand*\DialogueLabel[1]{%
  {#1}:%
}

% B. Thiết lập với độ dài tên dài nhất. Mặc định chọn Hanamaru từ Love Live!.
\newlength{\widestname}
\setlength{\widestname}{2.5cm}

% C. Tái định nghĩa lại hàm dialogue
\makeatletter
\renewenvironment{dialogue} {%
    \begin{list}{} {%
        \setlength\itemsep{\z@ \@plus .5ex}%
        \setlength{\parsep}{\parskip}%
        \setlength{\rightmargin}{0pt}           % Không căn lề phải
        \setlength{\labelwidth}{\widestname}    % Thiết đặt khoảng cách với tên dài nhất
        \setlength{\labelsep}{0.5em}            % Thiết đặt khoảng cách giữa tên nhân vật và đoạn hội thoại
        \setlength{\leftmargin}{\labelwidth}    % Căn lề trái
        \addtolength{\leftmargin}{\labelsep}    % 
        \defcommand\speak [1] {\item[{##1}]}    % Thiết lập lệnh speak trong dialogue
        \let\makelabel\DialogueLabel
      }%
      \PreDialogue\relax
    }{%
  \end{list}%
  }
\makeatother

% D. Hướng mũi tên
\newcommand{\arrow}[1][-45]
      {\begin{tikzpicture}[baseline = -0.5ex]
          \node[inner sep=0pt,outer sep=0pt,rotate = #1] (a) at (0,0)  {\rightarrow};
       \end{tikzpicture}}

% E. Chiêu mộ toàn bộ ký tự
\makeatletter
\font\SOUL@tt="Times New Roman"
\setbox\z@\hbox{\SOUL@tt-}
\SOUL@ttwidth\wd\z@ %reset default width of -
\makeatother

% F. Tabbing
\newcommand\tabs[1][1cm]{\hspace*{#1}}

% G. Tạo môi trường đề căn lề trái-phải
\newcommand{\tl}[1]{\noindent\begin{flushleft} #1 \end{flushleft}}
\newcommand{\tc}[1]{\noindent\begin{center} #1 \end{center}}
\newcommand{\tr}[1]{\noindent\begin{flushright} #1 \end{flushright}}

%=================================ĐẦU TRANG VÀ CHÂN TRANG=====================================
\pagestyle{fancy}
\fancyhead{} 
% Tên cuốn truyện
\fancyhead[LO]
{\device\fontsize{18pt}{18pt}\selectfont
\color{vol}{The Aogami High's Curse}
\\[-18pt]}                      
\fancyhead[RE]
{\device\fontsize{18pt}{18pt}\selectfont
\color{vol}{La malédiction du lycée Aogami}
\\[-18pt]}                     
% Số chương
\fancyhead[RO]
{\sen\fontsize{18pt}{18pt}\selectfont
\color{vols}{error}
\ab\fontsize{18pt}{18pt}\selectfont
\color{vols}{\textbf{(ER)}}
\\[-18pt]
}
\fancyhead[LE]
{\sen\fontsize{18pt}{18pt}\selectfont
\color{vols}{\textbf{erreur}}
\ab\fontsize{18pt}{18pt}\selectfont
\color{vols}{\textbf{(ER)}}
\\[-18pt]}
% Chân trang
\fancyfoot{} 
\fancyfoot[CO, CE]{\thepage}
\fancyfoot[LO]
{\ab\fontsize{10pt}{10pt}\selectfont \color{vol} [DIGITAL FORMAT ONLY ⇒ DO NOT PRINT]}
\fancyfoot[RE]
{\ab\fontsize{10pt}{10pt}\selectfont \color{vol} [FORMAT NUMÉRIQUE UNIQUEMENT ⇒ NE PAS IMPRIMER]}

%==========================================PHẦN THÂN==========================================
\begin{document}

%===========================================MỤC LỤC===========================================
% Trang đầu tiên
\newpage
\thispagestyle{empty}
\

\vspace*{\fill}
\noindent
\begin{minipage}[c][0.25\textheight][b]{\textwidth}
    \centering
    \includegraphics[width=15cm]{../../../Logo/round1.png}
    \vspace{30pt} % Khoảng cách từ logo xuống tâm trang
\end{minipage}\par
\nointerlineskip
\noindent
\begin{minipage}[c][0.25\textheight][t]{\textwidth}
    \vspace{30pt} % Khoảng cách từ tâm trang xuống phần chữ
    \centering
    {\huge \color{vols}{{\frank volume \hbrk ERROR}}} \vspace{15pt}

    {\fontsize{50pt}{50pt}\selectfont \color{vol}{\vnmdisp (LỜI NGUYỀN Ở\\THỊ TRẤN AOGAMI)}}\vspace{15pt}

    % {\Large \color{vol}{\vnmsans (Mùa thu năm 2021)}}
\end{minipage}
\vspace*{\fill}

\thispagestyle{empty}
\renewcommand{\contentsname}
{\centering\vnmsans\Large\textcolor{vol}{Mục lục}}
{
    \hypersetup{linkcolor=black}
    \tableofcontents
}
%========================================LỜI NÓI ĐẦU==========================================
\newpage
\section*{\phantomsection\label{preface}}
\addcontentsline{toc}{section}{Giới thiệu}
\vspace{-30pt}
{\centering\vnmsans\Large\textcolor{vol}{Giới thiệu}\par}
\subfile{intro}

\newpage
\section*{\phantomsection\label{nametable}}
\addcontentsline{toc}{section}{Bảng quy đổi tên các thành viên từ hololive gốc sang hololive ERROR}
\vspace{-30pt}
{\centering\vnmsans\Large\textcolor{vol}{Bảng quy đổi tên các thành viên từ hololive gốc sang hololive ERROR}\par}
\subfile{nametable}\

%=======================================DANH SÁCH TẬP========================================
%============================Phần 0: Mơ hồ============================
% \newpage
% \thispagestyle{empty}
% \begin{flushleft}
%         \hbox
%         {%
%             \makebox[\dimtorightedge]{}%
%             \makebox[0pt][r]
%             {\raisebox{0pt}[\dimtotop]{\includegraphics[width=\paperwidth]{../../../Image/BookER/chaptitle/chap0.jpg}}}%
%         }
% \end{flushleft}
% \vspace{-100pt}
% \section*{\phantomsection\label{ch.0}} 
% \addcontentsline{toc}{section}{\vnmsans{Mở đầu: Mơ hồ - 「\ruby{漠}{ばく}\ruby{然}{ぜん}」}} 
% \subsection*{\phantomsection\label{epE.P}}
% \addcontentsline{toc}{{subsection}{\flaregothic{ÉPISODE \frank 0}} - Tiếng gọi từ vòng lặp}
% \subfile{Episode/ER0}

%============================Phần 1: Quá khứ & Hiện đại============================
% \newpage
% \thispagestyle{empty}
% \begin{flushleft}
%         \hbox
%         {%
%             \makebox[\dimtorightedge]{}%
%             \makebox[0pt][r]
%             {\raisebox{0pt}[\dimtotop]{\includegraphics[width=\paperwidth]{../../../Image/BookER/chaptitle/chap1.jpg}}}%
%         }
% \end{flushleft}
% \vspace{-50pt}
% \thispagestyle{empty}
% \subfile{Episode/introch1}

% \newpage
% \section*{\phantomsection\label{ch.1}} 
% \addcontentsline{toc}{section}{\vnmsans{Phần Một: Hiện đại và Quá khứ - 「\ruby{現}{げん}\ruby{代}{だい}・\ruby{過}{か}\ruby{去}{こ}」}} 
% \subsection*{\phantomsection\label{epE1.1}} 
% \addcontentsline{toc}{{subsection}{\flaregothic{ÉPISODE \frank 1}} - Chào mừng tới Trường Trung học Aogami!}
% \vspace{-60pt}
% \subfile{Episode/ER1-01}

% \newpage
% \subsection*{\phantomsection\label{epE1.2}} 
% \addcontentsline{toc}{{subsection}{\flaregothic{ÉPISODE \frank 2}} - Lời đồn về học sinh chuyển trường}
% \vspace{-30pt}
% \subfile{Episode/ER1-02}

% \newpage
% \subsection*{\phantomsection\label{epE1.3}} 
% \addcontentsline{toc}{{subsection}{\flaregothic{ÉPISODE \frank 3}} - Sự sụp đổ của Trường Trung học Aogami}
% \vspace{-30pt}
% \subfile{Episode/ER1-03}

% \newpage
% \subsection*{\phantomsection\label{epE1.4}} 
% \addcontentsline{toc}{{subsection}{\flaregothic{ÉPISODE \frank 4}} - Ta cùng về nhà thôi, những học sinh chuyển trường}
% \vspace{-30pt}
% \subfile{Episode/ER1-04}

% \newpage
% \subsection*{\phantomsection\label{epE1.5}} 
% \addcontentsline{toc}{{subsection}{\flaregothic{ÉPISODE \frank 5}} - Lời ngiuyền của Misora Shino}
% \vspace{-30pt}
% \subfile{Episode/ER1-05}

% \newpage
% \subsection*{\phantomsection\label{epE1.6}} 
% \addcontentsline{toc}{{subsection}{\flaregothic{ÉPISODE \frank 6}} - Kẻ noi gương}
% \vspace{-30pt}
% \subfile{Episode/ER1-06}

% \newpage
% \subsection*{\phantomsection\label{epE1.7}} 
% \addcontentsline{toc}{{subsection}{\flaregothic{ÉPISODE \frank 7}} - Đêm thử thách lòng can đảm định mệnh}
% \vspace{-30pt}
% \subfile{Episode/ER1-07}

% \newpage
% \subsection*{\phantomsection\label{epE1.8}} 
% \addcontentsline{toc}{{subsection}{\flaregothic{ÉPISODE \frank 8}} - Lời ước nguyện dang dở}
% \vspace{-30pt}
% \subfile{Episode/ER1-08}

%============================Phần 2: Tàu điện============================
% \newpage
% \thispagestyle{empty}
% \begin{flushleft}
%         \hbox
%         {%
%             \makebox[\dimtorightedge]{}%
%             \makebox[0pt][r]
%             {\raisebox{0pt}[\dimtotop]{\includegraphics[width=\paperwidth]{../../../Image/BookER/chaptitle/chap2.jpg}}}%
%         }
% \end{flushleft}
% \vspace{-50pt}
% \thispagestyle{empty}
% \subfile{Episode/introch2}

% \newpage
% \section*{\phantomsection\label{ch.2}} 
% \addcontentsline{toc}{section}{\vnmsans{Phần Hai: Tàu điện - 「\ruby{電}{でん}\ruby{車}{しゃ}」}}
% \subsection*{\phantomsection\label{epE2.1}} 
% \addcontentsline{toc}{{subsection}{\flaregothic{ÉPISODE \frank 9}} - Bóng ma trong nhà ga Aogami}
% \vspace{-60pt}
% \subfile{Episode/ER2-01}

% \newpage
% \subsection*{\phantomsection\label{epE2.2}} 
% \addcontentsline{toc}{{subsection}{\flaregothic{ÉPISODE \frank 10}} - Sự bất thường trong sự bình thường}
% \vspace{-30pt}
% \subfile{Episode/ER2-02}

% \newpage
% \subsection*{\phantomsection\label{epE2.3}} 
% \addcontentsline{toc}{{subsection}{\flaregothic{ÉPISODE \frank 11}} - Con tàu ma}
% \vspace{-30pt}
% \subfile{Episode/ER2-03}

% \newpage
% \subsection*{\phantomsection\label{epE2.4}} 
% \addcontentsline{toc}{{subsection}{\flaregothic{ÉPISODE \frank 12}} - Mất tích}
% \vspace{-30pt}
% \subfile{Episode/ER2-04}

% \newpage
% \subsection*{\phantomsection\label{epE2.5}} 
% \addcontentsline{toc}{{subsection}{\flaregothic{ÉPISODE \frank 13}} - Yuka đã đi đâu?}
% \vspace{-30pt}
% \subfile{Episode/ER2-05}

% \newpage
% \subsection*{\phantomsection\label{epE2.6}} 
% \addcontentsline{toc}{{subsection}{\flaregothic{ÉPISODE \frank 14}} - Ký ức bị lãng quên}
% \vspace{-30pt}
% \subfile{Episode/ER2-06}

% \newpage
% \subsection*{\phantomsection\label{epE2.7}} 
% \addcontentsline{toc}{{subsection}{\flaregothic{ÉPISODE \frank 15}} - Tại sao\dots?}
% \vspace{-30pt}
% \subfile{Episode/ER2-07}

% \newpage
% \subsection*{\phantomsection\label{epE2.8}} 
% \addcontentsline{toc}{{subsection}{\flaregothic{ÉPISODE \frank 16}} - Tôi\dots\ là ai?}
% \vspace{-30pt}
% \subfile{Episode/ER2-08}

%============================Phần 3: Sự thật============================
% \newpage
% \thispagestyle{empty}
% \begin{flushleft}
%         \hbox
%         {%
%             \makebox[\dimtorightedge]{}%
%             \makebox[0pt][r]
%             {\raisebox{0pt}[\dimtotop]{\includegraphics[width=\paperwidth]{../../../Image/BookER/chaptitle/chap3.jpg}}}%
%         }
% \end{flushleft}
% \vspace{-50pt}
% \thispagestyle{empty}
% \subfile{Episode/introch3}

% \newpage
% \section*{\phantomsection\label{ch.3}} 
% \addcontentsline{toc}{section}{\vnmsans{Phần Ba: Sự thật - 「\ruby{真}{しん}\ruby{相}{そう}」}}
% \subsection*{\phantomsection\label{epE3.1}} 
% \addcontentsline{toc}{{subsection}{\flaregothic{ÉPISODE \frank 17}} - Chấp nhận thực tại}
% \vspace{-60pt}
% \subfile{Episode/ER3-01}

% \newpage
% \subsection*{\phantomsection\label{epE3.2}} 
% \addcontentsline{toc}{{subsection}{\flaregothic{ÉPISODE \frank 18}} - Ký ức từ khu rừng}
% \vspace{-30pt}
% \subfile{Episode/ER3-02}

% \newpage
% \subsection*{\phantomsection\label{epE3.3}} 
% \addcontentsline{toc}{{subsection}{\flaregothic{ÉPISODE \frank 19}} - Sự thật nghiệt ngã}
% \vspace{-30pt}
% \subfile{Episode/ER3-03}

% \newpage
% \subsection*{\phantomsection\label{epE3.4}} 
% \addcontentsline{toc}{{subsection}{\flaregothic{ÉPISODE \frank 20}} - Thế giới hạnh phúc giả tạo}
% \vspace{-30pt}
% \subfile{Episode/ER3-04}
    
% \newpage
% \subsection*{\phantomsection\label{epE3.5}} 
% \addcontentsline{toc}{{subsection}{\flaregothic{ÉPISODE \frank 21}} - Những ngày trong niềm hối hận}
% \vspace{-30pt}
% \subfile{Episode/ER3-05}

% \newpage
% \subsection*{\phantomsection\label{epE3.6}} 
% \addcontentsline{toc}{{subsection}{\flaregothic{ÉPISODE \frank 22}} - Lời ai oán của Sakura}
% \vspace{-30pt}
% \subfile{Episode/ER3-06}

% \newpage
% \subsection*{\phantomsection\label{epE3.7}} 
% \addcontentsline{toc}{{subsection}{\flaregothic{ÉPISODE \frank 23}} - Lời mong ước của Shino}
% \vspace{-30pt}
% \subfile{Episode/ER3-07}

% \newpage
% \subsection*{\phantomsection\label{epE3.8}} 
% \addcontentsline{toc}{{subsection}{\flaregothic{ÉPISODE \frank 24}} - Chào mừng tới Trường Trung học Aogami, một lần nữa!}
% \vspace{-30pt}
% \subfile{Episode/ER3-08}

%============================Phần 4: Khôi phục============================
\newpage
\thispagestyle{empty}
\begin{flushleft}
    \hbox
    {%
        \makebox[\dimtorightedge]{}%
        \makebox[0pt][r]
        {\raisebox{0pt}[\dimtotop]{\includegraphics[width=\paperwidth]{../../../Image/BookER/chaptitle/chap4.jpg}}}%
    }
\end{flushleft}
\vspace{-50pt}
\thispagestyle{empty}
\subfile{Episode/introch4}

\newpage
\section*{\phantomsection\label{ch.4}}
\addcontentsline{toc}{section}{\vnmsans{Phần Ba: Khôi phục - 「\ruby{回}{かい}\ruby{復}{ふく}」}}
\subsection*{\phantomsection\label{epEX.1}}
\addcontentsline{toc}{{subsection}{\flaregothic{ÉPISODE \frank 25}} - Hồn ma trẻ em}
\vspace{-60pt}
\subfile{Episode/ERX-01}

\newpage
\subsection*{\phantomsection\label{epEX.2}}
\addcontentsline{toc}{{subsection}{\flaregothic{ÉPISODE \frank 26}} - Cuộc sống học đường thường nhật}
\vspace{-30pt}
\subfile{Episode/ERX-02}

\newpage
\subsection*{\phantomsection\label{epEX.3}}
\addcontentsline{toc}{{subsection}{\flaregothic{ÉPISODE \frank 27}} - Bữa tiệc chào đón học sinh mới}
\vspace{-30pt}
\subfile{Episode/ERX-03}

\newpage
\subsection*{\phantomsection\label{epEX.4}}
\addcontentsline{toc}{{subsection}{\flaregothic{ÉPISODE \frank 28}} - Ngày cuối tuần hoạt náo}
\vspace{-30pt}
\subfile{Episode/ERX-04}

\newpage
\subsection*{\phantomsection\label{epEX.5}}
\addcontentsline{toc}{{subsection}{\flaregothic{ÉPISODE \frank 29}} - Phiêu lưu trong rừng ma ám}
\vspace{-30pt}
\subfile{Episode/ERX-05}

\newpage
\subsection*{\phantomsection\label{epEX.6}}
\addcontentsline{toc}{{subsection}{\flaregothic{ÉPISODE \frank 30}} - Hội thao trường học}
\vspace{-30pt}
\subfile{Episode/ERX-06}

\newpage
\subsection*{\phantomsection\label{epEX.7}}
\addcontentsline{toc}{{subsection}{\flaregothic{ÉPISODE \frank 31}} - Tiếng vọng quá khứ}
\vspace{-30pt}
\subfile{Episode/ERX-07}

\newpage
\subsection*{\phantomsection\label{epEX.8}}
\addcontentsline{toc}{{subsection}{\flaregothic{ÉPISODE \frank 32}} - Liều thuốc tình yêu}
\vspace{-30pt}
\subfile{Episode/ERX-08}

\newpage
\subsection*{\phantomsection\label{epEX.9}}
\addcontentsline{toc}{{subsection}{\flaregothic{ÉPISODE \frank 33}} - Mảnh ký ức cuối cùng}
\vspace{-30pt}
\subfile{Episode/ERX-09}

\newpage
\subsection*{\phantomsection\label{epEX.10}}
\addcontentsline{toc}{{subsection}{\flaregothic{ÉPISODE \frank 34}} - Sự bình yên sau cơn bão}
\vspace{-30pt}
\subfile{Episode/ERX-10}

\newpage
\subsection*{\phantomsection\label{epEX.11}}
\addcontentsline{toc}{{subsection}{\flaregothic{ÉPISODE \frank 35}} - Những mảnh ghép của ánh sáng}
\vspace{-30pt}
\subfile{Episode/ERX-11}

\newpage
\subsection*{\phantomsection\label{epEX.12}}
\addcontentsline{toc}{{subsection}{\flaregothic{ÉPISODE \frank 36}} - Cuộc chiến tri thức}
\vspace{-30pt}
\subfile{Episode/ERX-12}

\newpage
\subsection*{\phantomsection\label{epEX.13}}
\addcontentsline{toc}{{subsection}{\flaregothic{ÉPISODE \frank 37}} - Ngày lễ hội cuồng nhiệt}
\vspace{-30pt}
\subfile{Episode/ERX-13}

\newpage
\subsection*{\phantomsection\label{epEX.14}}
\addcontentsline{toc}{{subsection}{\flaregothic{ÉPISODE \frank 38}} - Hẹn gặp lại, Aogami!}
\vspace{-30pt}
\subfile{Episode/ERX-14}
%===========================================LỜI KẾT===========================================
\newpage
\section*{\phantomsection\label{end}}
\addcontentsline{toc}{section}{Lời kết}
{\centering\vnmsans\Large\textcolor{vol}{Lời kết}\par}
\subfile{end}

% \vfill % Đẩy phần dưới về đáy trang
% \begin{center}
%     \textcolor{vol}
%     {\quincy\Huge\textbf{TOME ERREUR TERMINE !} \\[1ex]
%         \vnmdisp\Huge\textbf{(QUYỂN ERROR HOÀN THÀNH!)}}
% \end{center}

\end{document}